\documentclass{article}
\usepackage[margin=1in]{geometry} % margins
\usepackage[justification=centering]{caption}
%\usepackage{calrsfs}
\usepackage{amsmath}
\usepackage{amsfonts}
\usepackage{amssymb}
\usepackage{parskip}
\usepackage{graphicx}
\usepackage{enumerate}
\usepackage{enumitem}
\usepackage{amsthm}
\usepackage{float}
\usepackage{tikz-cd}
\usepackage{ytableau}

\renewcommand{\implies}{\Rightarrow}
\newcommand{\longimplies}{\Longrightarrow}
\newcommand{\N}{\mathbb{N}}
\newcommand{\Z}{\mathbb{Z}}
\newcommand{\Q}{\mathbb{Q}}
\newcommand{\R}{\mathbb{R}}
\newcommand{\C}{\mathbb{C}}
\newcommand{\D}{\mathbb{D}}
\renewcommand{\H}{\mathbb{H}}
\newcommand{\F}{\mathbb{F}}
\newcommand{\K}{\mathbb{K}}
\newcommand{\Lf}{\mathbb{L}}
\newcommand{\Sp}{\mathbb{S}}
\newcommand{\cl}[1]{\overline{#1}}
\renewcommand{\emptyset}{\varnothing}
\newcommand{\norm}[1]{\left\|#1\right\|}
\newcommand{\maxnorm}[1]{\left\|#1\right\|_{\text{max}}}
\newcommand{\opnorm}[1]{\left\|#1\right\|_{\mathrm{op}}}
\newcommand{\supnorm}[1]{\left\|#1\right\|_{\infty}}
\newcommand{\id}{\operatorname{id}}
\newcommand{\sspan}{\operatorname{span}}
\newcommand{\tr}{\operatorname{tr}}
\newcommand{\ch}{\operatorname{char}}
\newcommand{\poly}{\mathrm{P}}
\newcommand{\mtx}{\mathcal{M}}
\newcommand{\seq}{\mathbf{Seq}}
\newcommand{\vv}[1]{\mathbf{#1}}
\newcommand{\dirsum}{\oplus}
\newcommand{\rank}{\operatorname{rank}}
\newcommand{\im}{\operatorname{im}}

\newcommand{\cj}[1]{\overline{#1}}
\newcommand{\inprod}[2]{\left\langle#1,#2\right\rangle}
\renewcommand{\Re}{\operatorname{Re}}
\renewcommand{\Im}{\operatorname{Im}}

\newcommand{\res}{\operatorname{res}}

\newcommand{\floor}[1]{\left\lfloor #1 \right\rfloor}

% Theorem environments
\usepackage{amsthm}

\theoremstyle{definition} % The "plain" style italicizes all body text.
	\newtheorem{thm}{Theorem}
		\numberwithin{thm}{section} % Theorem numbers are determined by section.
	\newtheorem{lemma}[thm]{Lemma}
	\newtheorem{prop}[thm]{Proposition}
	\newtheorem{cor}[thm]{Corollary}
	\newtheorem{conj}[thm]{Conjecture}

\theoremstyle{definition} % The "definition" style does not italicize body text..
    \newtheorem{defn}[thm]{Definition}
	\newtheorem{example}[thm]{Example}
    \newtheorem{exercise}[thm]{Exercise}

\usepackage{mathtools}

\usetikzlibrary{decorations.markings,positioning,decorations.text}

\def\xr{5}
\def\yr{5}

\usepackage{titlesec}
\titleformat{\section}[hang]{\normalfont\fontsize{14}{0}\bfseries}{Problem 8.\thesection\ }{0pt}{}

\begin{document}

\section{}

Let
\[
\hat{f}(n) := \frac{1}{2\pi} \int_{-\pi}^{\pi} e^{-inx} f(x) \,dx
\]
so that the Fourier series of a function $f$ is
\[
\sum_{n \in \mathbb{Z}} \hat{f}(n) e^{inx}.
\]

    Show that
\[
\sum_{j=-N}^{N} e^{ijx} = \frac{\sin((N+1/2)x)}{\sin(x/2)}.
\]
\begin{proof}
Let $f(x)$ be the function on the RHS. Notice that for any $x$ on which
$f(x)$ is defined, we have
\[
    f(x + 2\pi) = \frac{\sin((N+1/2)x + (2N + 1)\pi)}{\sin(x/2 + \pi)} 
    = \frac{-\sin((N+1/2)x )}{-\sin(x/2)} = f(x)
\]
that is, $f$ is $2\pi$-periodic where it is defined. In its current form,
$f$ is not defined at $2n\pi$ for $n \in \Z$, however, we have
\[
    \lim_{x \to 2n\pi} f(x) = \lim_{x \to 2n\pi} 
    \frac{\sin((N+1/2)x)}{\sin(x/2)}
    = \lim_{x \to 2n\pi} 
    \frac{(N + 1/2)\cos((N+1/2)x)}{(1/2)\cos(x/2)}
    = 2N + 1
\]
so $f$ can be extended to a $2\pi$-periodic function on $\R$. 
This means that the Fourier series of $f$ on $[-\pi, \pi]$ is
valid on all of $\R$. We now compute that series. For any $n \in \Z$,
we have
\begin{align*}
    e^{-inx}f(x)
    &= e^{-inx}\frac{\sin((N+1/2)x)}{\sin(x/2)} \\
    &= e^{-inx}\frac{e^{i(N+1/2)x} - e^{-i(N+1/2)x}}
    {e^{ix/2} - e^{-ix/2}} \\
    &= e^{-inx}e^{-iNx}\frac{e^{i(2N + 1)x} - 1}{e^{ix} - 1} \\
    &= e^{-inx}e^{-iNx}\sum_{m = 0}^{2N} e^{imx} \\
    &= \sum_{m = -N}^{N} e^{i(m - n)x}
\end{align*}
For $m \neq n$, the term $e^{i(m - n)x}$ has an antiderivative 
$e^{i(m - n)x}/(i(m - n))$, which evaluates to 0 over the bounds
$-\pi$ to $\pi$. Thus the $n$-th term of the Fourier series of $f$ is
\[
    \hat{f}(n) = \begin{cases}
        \frac{1}{2\pi}\int_{-\pi}^\pi dx & 
        \text{there is some } |m| \leq N, m = n \\
        0 & \text{otherwise}
    \end{cases} = \begin{cases}
        1 & |n| \leq N \\
        0 & |n| \geq N
    \end{cases}
\]
which corresponds to the Fourier series
\[
    f(x) = \sum_{n \in \Z} \hat{f}(n)e^{inx} = \sum_{n = -N}^N e^{inx}  
\]
Therefore, proven.
\end{proof}

% QUESTION 2
\section{}

\begin{enumerate}[label=(\alph*)]
    \item Let $f(\theta) = |\theta|$ on $[-\pi, \pi]$.
    \item Show that
\[
\hat{f}(n) = \begin{cases}
\frac{\pi}{2}, & n = 0 \\
\frac{-1 + (-1)^n}{\pi n^2}, & n \neq 0.
\end{cases}
\]

    \item Since $\sum_{n} |\hat{f}(n)| < \infty$, we will see in class that
\[
f(\theta) = \sum_{n \in \mathbb{Z}} \hat{f}(n) e^{in\theta}.
\]
    \item Write the Fourier series of $f$ in terms of sine and cosine.

    \item By taking $\theta = 0$, show that
\[
\sum_{n=0}^{\infty} \frac{1}{(2n+1)^2} = \frac{\pi^2}{8},
\]
\[
\sum_{n=1}^{\infty} \frac{1}{n^2} = \frac{\pi^2}{6}.
\]
Note that
\[
\sum_{n=1}^{\infty} \frac{1}{(2n)^2} = \frac{1}{4} \sum_{n=1}^{\infty} \frac{1}{n^2}.
\]
\end{enumerate}

\begin{proof}
\textit{(b).} For $n = 0$, we have
\[
    \hat{f}(0) = \frac{1}{2\pi}\int_{-\pi}^\pi |x|e^{-i0x}dx
    = \frac{1}{2\pi}\int_{-\pi}^\pi |x|dx
    = \frac{\pi}{2}
\]
and for $n \neq 0$, we have
\begin{align*}
    \hat{f}(n)
    &= \frac{1}{2\pi}\int_{-\pi}^\pi |x|e^{-inx}dx \\
    &= \frac{1}{2\pi}\left(\int_{0}^\pi xe^{-inx}dx - 
    \int_{-\pi}^0 xe^{-inx}dx\right) \\
    &= \frac{1}{2\pi}\left(\int_{0}^\pi xe^{-inx}dx + 
    \int_0^\pi xe^{inx}dx\right) \\
    &= \frac{1}{\pi}\int_0^\pi x\cos(nx) dx \\
    &= \frac{1}{\pi}\left(\left.\frac{x\sin(nx)}{n}\right|_0^\pi 
    - \int_0^\pi \frac{\sin(nx)}{n}dx\right) \\
    &= \frac{1}{\pi}\left(\left.\frac{\cos(nx)}{n^2}\right|_0^\pi\right) \\
    &= \frac{-1 + (-1)^n}{\pi n^2}
\end{align*}
Thus the formula for $\hat{f}(n)$ is proven.

\textit{(d).} By (c), we have
\begin{align*}
    f(\theta) 
    &= \sum_{n \in \mathbb{Z}} \hat{f}(n) e^{in\theta} \\
    &= \sum_{n \in \mathbb{Z}} \hat{f}(n) \cos n\theta +
    i\sum_{n \in \mathbb{Z}} \hat{f}(n)\sin n\theta \\
    &= \frac{\pi}{2} + \sum_{n \neq 0} 
    \frac{-1 + (-1)^n}{\pi n^2}\cos n\theta +
    i\sum_{n \neq 0} \frac{-1 + (-1)^n}{\pi n^2}\sin n\theta \\
    &= \frac{\pi}{2} 
    + 2\sum_{n = 1}^\infty \frac{-1 + (-1)^n}{\pi n^2}\cos n\theta
\end{align*}

\textit{(e).} Substituting $\theta = 0$ to the above, we get
\[
    \sum_{n = 1}^\infty \frac{1 - (-1)^n}{n^2} = \frac{\pi^2}{4}
\]
The terms of the series above for even $n$ are zero, so this just becomes
\[
    \sum_{n = 1}^\infty \frac{1}{(2n + 1)^2} = \frac{\pi^2}{8}
\]
We now decompose the absolutely convergent series 
$\sum_{n = 1}^\infty 1/n^2$ in the following way
\[
    \sum_{n = 1}^\infty \frac{1}{n^2} = 
    \sum_{n = 1}^\infty \frac{1}{(2n)^2} 
    + \sum_{n = 1}^\infty \frac{1}{(2n + 1)^2} 
    = \frac{1}{4}\sum_{n = 1}^\infty \frac{1}{n^2} 
    + \sum_{n = 1}^\infty \frac{1}{(2n + 1)^2}
\]
We thus have the following:
\[
    \frac{3}{4}\sum_{n = 1}^\infty \frac{1}{n^2} 
    = \sum_{n = 1}^\infty \frac{1}{(2n + 1)^2} = \frac{\pi^2}{8}
    \longimplies \sum_{n = 1}^\infty \frac{1}{n^2} = \frac{\pi^2}{6}
\]
\end{proof}

% QUESTION 3
\section{}

Let $F_N(x)$ be the Fej\'er kernel from class,
\[
F_N(x) = \frac{1}{N} \sum_{n=0}^{N-1} D_n(x), \quad D_N(x) = \sum_{n=-N}^{N} e^{inx}.
\]
\begin{enumerate}[label=(\alph*)]
    \item Show that
\[
F_N(x) = \frac{1}{N} \frac{\sin^2(Nx/2)}{\sin^2(x/2)}.
\]

\item Show that
\[
\frac{1}{2\pi} \int_{-\pi}^{\pi} F_N(x) \,dx = 1,
\]
and that
\[
\int_{|x| > \delta} F_N(x) \,dx \to 0 \quad \text{as } N \to \infty.
\]
That is, $F_N(x)$ is an approximate identity in the sense discussed in class.

\item Suppose that $f: [-\pi, \pi] \to \mathbb{R}$ is continuous except at a point $\theta_0 \in (-\pi, \pi)$, where it has left and right limits $f(\theta_0^-)$ and $f(\theta_0^+)$ respectively. Show that
\[
\lim_{N \to \infty} (f * F_N)(\theta_0) = \frac{f(\theta_0^+) + f(\theta_0^-)}{2},
\]
where $(f * g)(x) := \frac{1}{2\pi} \int_{-\pi}^{\pi} f(x - y) g(y) \,dy$.
\end{enumerate}

\begin{proof}
\textit{(a).} Note that $D_n(x) = \sin((n + 1/2)x)/\sin(x/2)$ by
problem 1. We then have the following:
\begin{align*}
    \sum_{n = 0}^{N - 1} D_n(x)
    &= \frac{1}{\sin(x/2)}\sum_{n = 0}^{N - 1} \sin((n + 1/2)x) \\
    &= \frac{1}{\sin(x/2)}
    \Im\left(\sum_{n = 0}^{N - 1} e^{i(n + 1/2)x}\right) \\
    &= \frac{1}{\sin(x/2)} \Im\left(e^{ix/2}
    \frac{e^{iNx} - 1}{e^{ix} - 1}\right) \\
    &= \frac{1}{\sin(x/2)} \Im\left(e^{iNx/2}
    \frac{e^{iNx/2} - e^{-iNx/2}}{e^{ix/2} - e^{-ix/2}}\right) \\
    &= \frac{1}{\sin(x/2)} \Im\left(e^{iNx/2}
    \frac{\sin(Nx/2)}{\sin(x/2)}\right) \\
    &= \frac{\sin^2(Nx/2)}{\sin^2(x/2)}
\end{align*}
Hence the formula for $F_N$ is proven.

\textit{(b).} For the first one, we have
\begin{align*}
    \frac{1}{2\pi} \int_{-\pi}^{\pi} F_N(x) \,dx
    &= \frac{1}{2N\pi}
    \sum_{n=0}^{N-1} \sum_{m=-n}^{n} \int_{-\pi}^{\pi} e^{imx} \,dx \\
    &= \frac{1}{2N\pi}\sum_{n=0}^{N-1} \left(
    \int_{-\pi}^{\pi} dx + \sum_{m=-n, m \neq 0}^{n} 
    \int_{-\pi}^{\pi} e^{imx} \,dx\right) \\
    &= \frac{1}{2N\pi}\sum_{n=0}^{N-1} \left(2\pi 
    + \sum_{m=-n, m \neq 0}^{n} 
    \left.\frac{e^{imx}}{im} \right|_{-\pi}^{\pi}\right) \\
    &= 1
\end{align*}

Let $0 < \delta < \pi$.
Notice that $F_N(x)$ is bounded above by $\frac{1}{N}\csc^2(x/2)$. $\csc$
is continuous on the compact subset $[-\pi, -\delta] \cup [\delta, \pi]$.
The upper bound $\frac{1}{N}\csc^2(x/2)$ then converges to 0 uniformly
so the integral (which is positive) 
\[
    \int_{|x| > \delta} F_N(x)dx \to \int_{|x| > \delta} 0dx = 0
\]

\textit{(c).} We first note that $F_N$ is an even function, so 
$\frac{1}{2\pi}\int_0^\pi F_N(x)dx = \frac{1}{2}$.
The convolution $(f * F_N)(\theta_0)$ is
\[
    (f * F_N)(\theta_0) 
    = \frac{1}{2\pi} \int_{-\pi}^{\pi} f(\theta_0 - y) F_N(y) \,dy
\]
Given $\varepsilon > 0$, we take $\delta > 0$ such that 
$|f(\theta_0 - y) - f(\theta_0^+)| < \varepsilon$ for 
$y \in (-\delta, 0)$ and 
$|f(\theta_0 - y) - f(\theta_0^-)| < \varepsilon$ for
$y \in (0, \delta)$. This prompts us to rewrite the convolution as
\[
    (f * F_N)(\theta_0) 
    = \frac{1}{2\pi} \int_{|y| \geq \delta} f(\theta_0 - y) F_N(y) \,dy
    + \frac{1}{2\pi} \int_0^{\delta} f(\theta_0 - y) F_N(y) \,dy
    + \frac{1}{2\pi} \int_{-\delta}^0 f(\theta_0 - y) F_N(y) \,dy
\]
Given a bound $M > 0$ for the function $f$, the first integral can be 
bounded between multiples of\\ $\int_{|x| > \delta} F_N(x)dx$ which
converges to 0 as $N \to \infty$. The second one can be bounded
in the following way:
\begin{align*}
    \frac{1}{2\pi} \int_0^{\delta} 
    (f(\theta_0^-) - \varepsilon) F_N(y) \,dy &\leq
    \frac{1}{2\pi} \int_0^{\delta} f(\theta_0 - y) F_N(y) \,dy
    \leq \frac{1}{2\pi} \int_0^{\delta} 
    (f(\theta_0^-) + \varepsilon) F_N(y) \,dy \\
    \longimplies
    (f(\theta_0^-) - \varepsilon)\frac{1}{2\pi} 
    \int_0^{\delta} F_N(y) \,dy &\leq
    \frac{1}{2\pi} \int_0^{\delta} f(\theta_0 - y) F_N(y) \,dy
    \leq (f(\theta_0^-) + \varepsilon)\frac{1}{2\pi} 
    \int_0^{\delta} F_N(y) \,dy \\
\end{align*}
As $N \to \infty$, $\frac{1}{2\pi}\int_0^{\delta} F_N(y) \,dy 
\to \frac{1}{2}$ as a consequence of (b). We thus have
\[
    \frac{f(\theta_0^-) - \varepsilon}{2} \leq
    \frac{1}{2\pi} \int_0^{\delta} f(\theta_0 - y) F_N(y) \,dy
    \leq \frac{f(\theta_0^-) + \varepsilon}{2} 
\]
Similarly, we get the following bound for the third integral:
\[
    \frac{f(\theta_0^+) - \varepsilon}{2} \leq
    \frac{1}{2\pi} \int_0^{\delta} f(\theta_0 - y) F_N(y) \,dy
    \leq \frac{f(\theta_0^+) + \varepsilon}{2} 
\]
We thus have that $\lim_{N \to \infty}(f * F_N)(\theta_0)$ is bounded
between $\frac{f(\theta_0^+) + f(\theta_0^-)}{2} - \varepsilon$
and $\frac{f(\theta_0^+) + f(\theta_0^-)}{2} + \varepsilon$, so
taking $\varepsilon \to 0$, we get
\[
    \lim_{N \to \infty} (f * F_N)(\theta_0) = \frac{f(\theta_0^+) + f(\theta_0^-)}{2}
\]
\end{proof}

% QUESTION 4
\section{}

Let $f: \mathbb{R} \to \mathbb{R}$ be $2\pi$-periodic such that $|f(x+h) - f(x)| \leq |h|^\alpha$ for all $x \in \mathbb{R}$ and $h \in \mathbb{R}$ for some $0 < \alpha < 1$. Show that
\[
|\hat{f}(n)| \leq C |n|^{-\alpha}
\]
for some $C > 0$.

\begin{proof}
For any $x \in \R$ and nonzero $k \in \Z$, we have
\begin{align*}
    \pi^\alpha|k|^{-\alpha}
    &\geq |f(x + \pi/k) -  f(x)| \\ 
    &= \left|\sum_{n \in \Z}\hat{f}(n)e^{in(x + \pi/k)}
    - \sum_{n \in \Z}\hat{f}(n)e^{inx}\right| \\
    &= \left|\sum_{n \in \Z}\hat{f}(n)e^{inx}
    \left(e^{i\pi n/k} - 1\right)\right| \\
\end{align*}
Note the following integral identity:
\[
    \int_{-\pi}^\pi e^{-ikx}\sum_{n \in \Z}\hat{f}(n)e^{inx}
    \left(e^{i\pi n/k} - 1\right)dx
    = \sum_{n \in \Z}\hat{f}(n)\left(e^{i\pi n/k} - 1\right)
    \int_{-\pi}^\pi e^{i(n - k)x}dx = -4\pi\hat{f}(k)
\]
We can now bound $|\hat{f}(k)|$:
\begin{align*}
    |\hat{f}(k)|
    &= \frac{4}{\pi}
    \left|\int_{-\pi}^\pi e^{-ikx}\sum_{n \in \Z}\hat{f}(n)e^{inx}
    \left(e^{i\pi n/k} - 1\right)dx\right| \\
    &\leq \frac{4}{\pi}
    \int_{-\pi}^\pi \left|e^{-ikx}\sum_{n \in \Z}\hat{f}(n)e^{inx}
    \left(e^{i\pi n/k} - 1\right)\right|dx \\
    &\leq 4\pi^{\alpha - 1}|k|^{-\alpha}
    \int_{-\pi}^\pi \left|e^{-ikx}\right|dx \\
    &= 4\pi^{\alpha - 1}|k|^{-\alpha}
    \int_{-\pi}^\pi dx \\
    &= 8\pi^{\alpha}|k|^{-\alpha}
\end{align*}
Therefore proven.
\end{proof}


% QUESTION 5
\section{}

Let
\[
V_n(x) := (1 + e^{inx} + e^{-inx}) F_n(x).
\]
\begin{enumerate}[label=(\alph*)]
\item Check that $\hat{V}_n(k) = 1$ for $|k| \leq n$. Show that $\|V_n'\|_1 \leq Cn$ for some $C > 0$ independent of $n$.

\item Suppose that $f$ is a trigonometric polynomial with $\hat{f}(k) = 0$ for $|k| > n$. Show that
\[
\|f'\|_1 \leq Cn \|f\|_1.
\]
Hint: Show that $f = V_n * f$, that $f' = V_n' * f$, and that $\|f * g\|_1 \leq \|f\|_1 \|g\|_1$ for any $f, g$.
\end{enumerate}

\begin{proof}
\textit{(a).} For $|k| \leq n$, we have
\begin{align*}
    \hat{V}_n(k)
    &= \frac{1}{2\pi}\int_{-\pi}^\pi 
    e^{-ikx}(1 + e^{inx} + e^{-inx}) F_n(x)dx \\
    &= \frac{1}{2n\pi}\sum_{m=0}^{n-1} 
    \sum_{j=-m}^{m}\int_{-\pi}^\pi 
    e^{-ikx}(1 + e^{inx} + e^{-inx}) e^{ijx}dx \\
    &= \frac{1}{2n\pi}\sum_{m=0}^{n-1} 
    \sum_{j=-m}^{m}\int_{-\pi}^\pi 
    (e^{-i(k - j)x} + e^{i(n + j - k)x} + e^{-i(n - j + k)x})dx \\
    &= \frac{1}{n}\sum_{m=0}^{n-1}\sum_{j=-m}^{m}
    (\chi[j = k] + \chi[j = k - n] + \chi[j = k + n])
\end{align*}
where $\chi[P]$, the indicator function, is 1 if $P$ holds and 0 otherwise. 
We now try evaluating that sum for each indicator separately:
\begin{align*}
    \sum_{m=0}^{n-1}\sum_{j=-m}^{m}\chi[j = k] 
    &= \sum_{m=|k|}^{n-1}\sum_{j=-m}^{m}\chi[j = k] 
    = \sum_{m=|k|}^{n-1} 1 
    = n - |k| \\
    \sum_{m=0}^{n-1}\sum_{j=-m}^{m}\chi[j = k - n]
    &=  \sum_{m=0}^{n-1}\sum_{j=n-m}^{n+m}\chi[j = k] \\
    &= \begin{cases}
        \sum_{m=n - k}^{n-1}\sum_{j=n-m}^{n+m}\chi[j = k] & k > 0\\
        0 & k \leq 0
    \end{cases} \\
    &= \begin{cases}
        k & k > 0\\
        0 & k \leq 0
    \end{cases} \\
    \sum_{m=0}^{n-1}\sum_{j=-m}^{m}\chi[j = k + n]
    &=  \sum_{m=0}^{n-1}\sum_{j=-n-m}^{-n+m}\chi[j = k] \\
    &= \begin{cases}
        \sum_{m=n + k}^{n-1}\sum_{j=-n-m}^{-n+m}\chi[j = k] & k < 0\\
        0 & k \geq 0
    \end{cases} \\
    &= \begin{cases}
        -k & k < 0\\
        0 & k \geq 0
    \end{cases}
\end{align*}
Adding all of these, we obtain
\[
    \sum_{m=0}^{n-1}\sum_{j=-m}^{m}
    (\chi[j = k] + \chi[j = k - n] + \chi[j = k + n]) = n
\]
and hence $\hat{V}_n(k) = 1$.

We write $V_n(x) = (1 + e^{inx} + e^{-inx})
\frac{\sin^2(nx/2)}{n\sin^2(x/2)}$ and compute its derivative:
\[
    V'_n(x) = (ine^{inx} - ine^{-inx})F_n(x) 
    + (1 + e^{inx} + e^{-inx}) F_n'(x)
\]
Thus we have 
\[
    \norm{V'_n}_1 \leq 2n\norm{F_n}_1 + 3\norm{F_n'}_1
\]
Since $F_n(x) \geq 0$ (by Question 2) and 
$\int_{-\pi}^\pi F_n(x)dx = 2\pi$, we have $\norm{F_n}_1 = 2\pi$.
For the norm of $F_n'$, we note that $F_n$ is a trigonometric 
polynomial with coefficients bounded by 1, we have that 
$|F_n'(x)| \leq n(n + 1) \leq An^2$ for some $A > 0$, where 
$A$ is independent of $n$. The derivative of $F_n$ is
\[
    F'_n(x) = \frac{\sin(nx/2)\cos(nx/2)}{\sin^2(x/2)} 
    - \frac{\sin^2(nx/2)\cos(x/2)}{n\sin^3(x/2)}
\]
Since $|\sin x| \leq |x|$ for all $x$, $|\sin x| \geq \pi|x|/2$ 
for all $|x| < \pi/2$, and sine and cosine is bounded above by 1, we have
\[
    |F'_n(x)| \leq \frac{1}{\pi^2x^2/4}
    + \frac{n|x|/2}{n\pi^3|x|^3/8} = \frac{4}{\pi^2x^2}
    + \frac{4}{\pi^3x^2} = \frac{B}{x^2}
\]
for $x \in [-\pi, \pi]$ and some constant $B > 0$ independent of $n$.
Then $\norm{F'_n}_1$ can be bounded in the following way
\begin{align*}
    \norm{F'_n}_1 = \int_{-\pi}^\pi |F'_n(x)|dx
    &= \int_{|x| \leq 1/n} |F'_n(x)|dx 
    + \int_{|x| \geq 1/n} |F'_n(x)|dx \\
    &\leq \int_{|x| \leq 1/n} An^2dx 
    + \int_{|x| \geq 1/n} \frac{B}{x^2}dx \\
    &= 2An + \frac{-2B}{x}\bigg|_{1/n}^\pi \\
    &= 2(A + B)n + \frac{-2B}{\pi}
\end{align*}
This means that for some constant $c > 0$ that depends only on 
$A$ and $B$ (i.e. $c$ is independent of $n$), we have 
$\norm{F'_n}_1 \leq cn$. We thus have the bound
\[
    \norm{V'_n}_1 \leq 2n\norm{F_n}_1 + 3\norm{F_n'}_1 \leq 
    4\pi n + 3cn = Cn
\]
for a constant $C > 0$ independent of $n$.

\textit{(b).}
Write $f(x) = \sum_{k = -n}^n \hat{f}(k)e^{ikx}$. To show that 
$f = V_n * f$, it suffices to show that their Fourier coefficients
are equal. For $|k| \leq n$, we have
\[
    \widehat{V_n * f}(k) = \hat{V_n}(k)\hat{f}(k) = \hat{f}(k)
\]
and for $|k| > n$,
\[
    \widehat{V_n * f}(k) = \hat{V_n}(k)\hat{f}(k) = 0 = \hat{f}(k)
\]
Thus $f = V_n * f$. Differentiating this equation, we get
\begin{align*}
    f'(x)
    &= \frac{d}{dx}\frac{1}{2\pi}\int_{-\pi}^\pi 
    V_n(x - y)f(y) dy \\
    &= \frac{1}{2\pi}\int_{-\pi}^\pi \frac{\partial}{\partial x}
    V_n(x - y)f(y) dy \\
    &= \frac{1}{2\pi}\int_{-\pi}^\pi 
    V_n'(x - y)f(y) dy \\
    &= V_n' * f(x)
\end{align*}
thus $f' = V_n' * f$. 
The 1-norm of the right hand side of this can be bounded
\begin{align*}
    \norm{V_n' * f}_1
    &= \int_{-\pi}^\pi |V_n' * f(x)|dx \\
    &= \int_{-\pi}^\pi \left|\frac{1}{2\pi}\int_{-\pi}^\pi 
    V_n'(x - y)f(y) dy\right|dx \\
    &\leq \frac{1}{2\pi}\int_{-\pi}^\pi \int_{-\pi}^\pi 
    \left|V_n'(x - y)f(y)\right|dydx \\
    &= \frac{1}{2\pi}\int_{-\pi}^\pi \int_{-\pi}^\pi 
    \left|V_n'(x - y)f(y)\right|dxdy \\
    &= \frac{1}{2\pi}\int_{-\pi}^\pi \int_{-\pi - y}^{\pi - y} 
    \left|V_n'(t)f(y)\right|dtdy \\
    &\leq \frac{1}{2\pi}\int_{-\pi}^\pi \int_{-\pi}^{\pi} 
    \left|V_n'(t)f(y)\right|dtdy \\
    &= \frac{1}{2\pi}\int_{-\pi}^\pi 
    \left|V_n'(t)\right|dt\int_{-\pi}^{\pi} 
    \left|f(y)\right|dy \\
    &\leq \norm{V_n'}_1\norm{f}_1
\end{align*}
By part (a), we obtain the following bound for the norm of $f'$:
\[
    \norm{f'}_1 \leq Cn\norm{f}_1
\]
\end{proof}

%QUESTION 6
\section{}

Continuing with notation from the previous problem, also for any $s > 0$ define the \textquotedblleft norm\textquotedblright (explain why this is not a norm):
\[
\|f\|_{H^s}^2 := |\hat{f}(0)|^2 + \sum_{n \in \mathbb{Z}} |n|^{2s} |\hat{f}(n)|^2.
\]
For $0 < \alpha < 1$ define
\[
\|f\|_{C^\alpha} := \sup_{x,h} \frac{|f(x+h) - f(x)|}{|h|^\alpha} + \sup_x |f(x)|.
\]

In this problem, we will prove a version of the Sobolev embedding theorem via a discrete version of Littlewood-Paley theory.
\begin{enumerate}[label=(\alph*)]
    \item  Let
\[
\varphi_j(x) = V_{2^{j-1}}(x)\left(e^{i 3 \cdot 2^{j-1} x} + e^{-i 3 \cdot 2^{j-1} x}\right).
\]
Check that $\hat{\varphi}_j(n) = 1$ for $2^j \leq |n| \leq 2^{j+1}$ and that $\hat{\varphi}(0) = 0$, which is the same as
\[
\int_{-\pi}^{\pi} \varphi_j(x) \,dx = 0.
\]

\item Show that
\[
|\varphi_j(x)| \leq C \frac{1}{2^j} \min(2^{2j}, |x|^{-2}).
\]
Hint: similar bounds hold for $F_n(x)$.

\item Show that
\[
\sum_{2^j \leq |n| \leq 2^{j+1}} |\hat{f}(n)|^2 \leq \int |(\varphi_j * f)(x)|^2 \,dx.
\]

\item Show that
\[
    |(\varphi_j * f)(x)| = \frac{1}{2\pi}\left|\int_{-\pi}^\pi
    \varphi_j(y)(f(x - y) - f(x))dy\right| 
    \leq C\norm{f}_{C^\alpha}2^{-\alpha j}
\]
Conclude that 
\[
    \sum_{2^j \leq |n| \leq 2^{j+1}} |\hat{f}(n)|^2 
    \leq C\norm{f}^2_{C^\alpha}2^{-2\alpha j}
\]

\item Let $0 < \alpha < 1$ and let $\beta < \alpha$. 
Show that there is a $C_{\alpha, \beta}$ s.t.
\[
    \norm{f}_{H^\beta} \leq C_{\alpha, \beta}\norm{f}_\alpha
\]
This is a version of the \textit{Sobolev embedding theorem}. 
The idea of splitting up the Fourier coefficients into dyadic 
blocks (of comparable frequencies) is known as
Littlewood-Paley theory and has many applications.

\item Conclude that if $f \in C^\alpha$ for $\alpha > \frac{1}{2}$, then
\[
\sum_n |\hat{f}(n)| < \infty.
\]
\end{enumerate}

\begin{proof}
\textit{(a).} Let $2^j \leq |n| \leq 2^{j + 1}$. We now compute the
Fourier coefficient $\hat{\varphi}_j(n)$:
\begin{align*}
    \hat{\varphi}_j(n)
    &= \frac{1}{2\pi}\int_{-\pi}^\pi e^{-inx}
    V_{2^{j-1}}(x)\left(e^{i 3 \cdot 2^{j-1} x} 
    + e^{-i 3 \cdot 2^{j-1} x}\right) dx \\
    &= \hat{V}_{2^{j-1}}\left(n - 3 \cdot 2^{j-1}\right)
    + \hat{V}_{2^{j-1}}\left(n + 3 \cdot 2^{j-1}\right)
\end{align*}
Note that $V_{2^{j - 1}}$ is a trigonometric polynomial of
degree $2^j - 1$
If $n > 0$ then $n + 3 \cdot 2^{j-1} \geq 2^j + 3 \cdot 2^{j-1} 
> 2^{j + 1}$, so the coefficient 
$\hat{V}_{2^{j-1}}\left(n + 3 \cdot 2^{j-1}\right) = 0$. On the 
other hand, we have $2^j \leq n \leq 2^{j + 1} \implies
-2^{j - 1} \leq n - 3 \cdot 2^{j - 1} \leq 2^{j - 1}$, so
by Question 5, the coefficient 
$\hat{V}_{2^{j-1}}\left(n + 3 \cdot 2^{j-1}\right) = 1$. Hence
$\hat{\varphi}_j(n) = 1$. The case $n < 0$ is handled almost identically.

For $n = 0$, we have
\[
    \hat{\varphi}_j(0) = \hat{V}_{2^{j-1}}\left(-3 \cdot 2^{j-1}\right)
    + \hat{V}_{2^{j-1}}\left(3 \cdot 2^{j-1}\right)
\]
and since $|3 \cdot 2^{j-1}| > 2^j$ and $V_{2^{j - 1}}$ is a 
trigonometric polynomial is of degree $2^{j - 1} - 1$.
Thus both of the coefficients in the sum above are 0.
Hence, $\hat{\varphi}_j(0) = 0$.

\textit{(b).}
Notice that $|\varphi_j(x)|$ can be bounded above in the following way:
\[
    |\varphi_j(x)| \leq 6|F_{2^{j - 1}}(x)|
\]
$F_{2^{j - 1}}$ is a trigonometric polynomial of degree $2^{j - 1} - 1$
with coefficients bounded by 1, so $|F_{2^{j - 1}}(x)|
\leq 2^j$ (remember to take negative terms into account), hence
\[
    |\varphi_j(x)| \leq 6\frac{1}{2^{j}}2^{2j}
\]
On the other hand, we have
\[
    |F_{2^{j - 1}}(x)| = \frac{1}{2^{j - 1}}
    \frac{\sin^2(2^{j - 1}x/2)}{\sin^2(x/2)} 
    \leq \frac{1}{2^{j - 1}}\frac{1}{x^2/\pi^2} 
    = \frac{\pi^2}{2^{j - 1}}|x|^{-2}
\]
Thus, we have
\[
    |\varphi_j(x)| \leq 12\pi^2\frac{1}{2^j}|x|^{-2}
\]
Hence, for some constant $C > 0$ (independent of $j$),
we have $|\varphi_j(x)| \leq C \frac{1}{2^j} \min(2^{2j}, |x|^{-2})$.

\textit{(c).} I'm speechless.
\begin{align*}
    \int_{-\pi}^\pi|\varphi_j * f(x)|^2 dx
    &\geq \frac{1}{2\pi}\int_{-\pi}^\pi|\varphi_j * f(x)|^2 dx \\
    &= \frac{1}{2\pi}\int_{-\pi}^\pi
    \left(\sum_n \hat{\varphi}_j(n)\hat{f}_j(n)e^{inx}\right)
    \cj{\left(\sum_m \hat{\varphi}_j(m)\hat{f}_j(m)e^{imx}\right)}dx \\
    &= \sum_{n,m}\frac{1}{2\pi}\int_{-\pi}^\pi
    \left(\hat{\varphi}_j(n)\hat{f}_j(n)\right)
    \cj{\left(\hat{\varphi}_j(m)\hat{f}_j(m)\right)}e^{i(n - m)x}dx \\
    &= \sum_{n}\left(\hat{\varphi}_j(n)\hat{f}_j(n)\right)
    \cj{\left(\hat{\varphi}_j(n)\hat{f}_j(n)\right)} \\
    &= \sum_{n} |\hat{\varphi}_j(n)|^2|\hat{f}_j(n)|^2
\end{align*}
We then have
\[
    \sum_{2^j \leq |n| \leq 2^{j+1}} |\hat{f}(n)|^2 
    = \sum_{2^j \leq |n| \leq 2^{j+1}} 
    |\hat{\varphi}_j(n)|^2|\hat{f}(n)|^2 
    \leq \sum_{n} 
    |\hat{\varphi}_j(n)|^2|\hat{f}(n)|^2 
    \leq \int_{-\pi}^\pi|\varphi_j * f(x)|^2 dx
\]

\textit{(d).} Still speechless
\begin{align*}
    |(\varphi_j * f)(x)| 
    &= \frac{1}{2\pi}\left|\int_{-\pi}^\pi
    \varphi_j(y)f(x - y)dy - f(x)\int_{-\pi}^\pi
    \varphi_j(y)dy\right| \\
    &= \frac{1}{2\pi}\left|\int_{-\pi}^\pi
    \varphi_j(y)(f(x - y) - f(x))dy\right| \\
    &\leq \frac{1}{2\pi}\int_{-\pi}^\pi
    |\varphi_j(y)|\cdot|y|^\alpha\norm{f}_\alpha dy \\
    &= \frac{1}{\pi}\int_{0}^{1/2^j}
    |\varphi_j(y)|\cdot|y|^\alpha\norm{f}_\alpha dy
    + \frac{1}{\pi}\int_{1/2^j}^\pi
    |\varphi_j(y)|\cdot|y|^\alpha\norm{f}_\alpha dy \\
    &\leq \frac{C\norm{f}_\alpha}{2^j\pi}\int_{0}^{1/2^j}
    2^{2j}|y|^\alpha dy
    + \frac{C\norm{f}_\alpha}{2^j\pi}\int_{1/2^j}^\pi
    |y|^{\alpha - 2} dy \\
    &= \frac{C\norm{f}_\alpha}{\pi}
    \frac{2^{-j\alpha}}{\alpha + 1}
    + \frac{C\norm{f}_\alpha}{2^j\pi}
    \left(\frac{\pi^{\alpha - 1}}{\alpha - 1}
    - \frac{2^{-j(\alpha - 1)}}{\alpha - 1}\right) \\
    &= \frac{C\norm{f}_\alpha}{\pi}
    \frac{2^{-j\alpha}}{\alpha + 1}
    - \frac{C\norm{f}_\alpha}{2^j\pi}\frac{\pi^{\alpha - 1}}{1 - \alpha}
    + \frac{C\norm{f}_\alpha}{\pi}\frac{2^{-j\alpha}}{1 - \alpha} \\
    &\leq C'\norm{f}_\alpha 2^{-j\alpha}
\end{align*}
Thus
\begin{align*}
    \sum_{2^j \leq |n| \leq 2^{j+1}} |\hat{f}(n)|^2
    &\leq \frac{1}{2\pi}\int_{-\pi}^\pi|\varphi_j * f(x)|^2 dx \\
    &\leq \frac{1}{2\pi}\int_{-\pi}^\pi
    C'^2\norm{f}^2_\alpha 2^{-2j\alpha} dx \\
    &= C\norm{f}^2_\alpha 2^{-2j\alpha}
\end{align*}

\textit{(e).}
I no want write words.
\begin{align*}
    \norm{f}^2_{H^\beta}
    &= |\hat{f}(0)|^2 + \sum_{n \neq 0}|n|^{2\beta}|\hat{f}(n)|^2 \\
    &\leq |\hat{f}(0)|^2 + 
    \sum_{j = 0}^\infty |2^{j + 1}|^{2\beta} 
    \sum_{2^j \leq |n| \leq 2^{j + 1}} |\hat{f}(n)|^2 \\
    &\leq |\hat{f}(0)|^2 + 
    2^{2\beta}\sum_{j = 0}^\infty 2^{2j\beta}
    C\norm{f}^2_\alpha 2^{-2j\alpha} \\
    &= |\hat{f}(0)|^2 + 
    2^{2\beta}C\norm{f}^2_\alpha
    \sum_{j = 0}^\infty \left(4^{\beta - \alpha}\right)^j \\
    &= |\hat{f}(0)|^2 + 
    2^{2\beta}C\norm{f}^2_\alpha \frac{1}{1 - 4^{\beta - \alpha}} \\
    &\leq C^2_{\alpha, \beta}\norm{f}^2_\alpha
\end{align*}
Thus
\[
    \norm{f}_{H^\beta} \leq C_{\alpha, \beta}\norm{f}_\alpha
\]

\textit{(f).}
Let $\alpha > 1/2$ and pick $1/2 < \beta < \alpha$.
Write $\sum_{n \neq 0} |\hat{f}(n)|$ as $\sum_{n \neq 0} |n|^{-\beta} 
\cdot |n|^{\beta}|\hat{f}(n)|$, we apply Holder's inequality to get
\[
    \left(\sum_{n \neq 0} |\hat{f}(n)|\right)^2
    \leq \left(\sum_{n \neq 0} |n|^{-2\beta}\right)
    \left(\sum_{n \neq 0} |n|^{2\beta}|\hat{f}(n)|^2\right)
    \leq \left(\sum_{n \neq 0} |n|^{-2\beta}\right)\norm{f}_{H^\beta}^2
    \leq \left(\sum_{n \neq 0} |n|^{-2\beta}\right)
    C_{\alpha, \beta}^2\norm{f}^2_\alpha
\]
The series in the latter expression converges since $2\beta > 1$, thus
$\sum_{n \neq 0} |\hat{f}(n)| < \infty$.
\end{proof}

\end{document}